\newpage
\section{第27章：案例研究——再论命令式对象}

\begin{taplauthorsolutionentrybox}
\phantomsection\label{sol:author:27.1}
\textbf{27.1 解答}

可以把 self 引用改为保存一个从表示记录到方法记录的函数，并让方法记录类型和
表示类型都成为有界类型参数：
\begin{lstlisting}
setCounterClass =
  lambda M<:SetCounter. lambda R<:CounterRep.
  lambda self:Ref (R->M). lambda r:R.
    {get = lambda _:Unit. !(r.x),
     set = lambda i:Nat. r.x:=i,
     inc = lambda _:Unit.
             (!self r).set (succ((!self r).get unit))};
? setCounterClass
  : All M<:SetCounter. All R<:CounterRep.
    (Ref (R->M)) -> R -> SetCounter

instrCounterClass =
  lambda M<:InstrCounter. lambda R<:InstrCounterRep.
  lambda self:Ref (R->M). lambda r:R.
    let super = setCounterClass [M] [R] self in
    {get = (super r).get,
     set = lambda i:Nat.
             (r.a:=succ(!(r.a)); (super r).set i),
     inc = (super r).inc,
     accesses = lambda _:Unit. !(r.a)};
? instrCounterClass
  : All M<:InstrCounter. All R<:InstrCounterRep.
    (Ref (R->M)) -> R -> InstrCounter

newInstrCounter =
  let m = ref
    (lambda r:InstrCounterRep. error as InstrCounter) in
  let m' = instrCounterClass
    [InstrCounter] [InstrCounterRep] m in
  (m := m';
   lambda _:Unit.
     let r = {x=ref 1, a=ref 0} in m' r);
? newInstrCounter : Unit -> InstrCounter
\end{lstlisting}
在父类实例化处，\(\tapltype{Ref}(R\to M)\) 两边完全相同，无须利用
\(\tapltype{Ref}\) 的协变性。

\taplsolutionbacklink{ex:27.1}{返回练习27.1}
\end{taplauthorsolutionentrybox}
